% The theory of splitting sets into two (inhabited, disjoint) partitions.

% See also "bisections" here: https://isarmathlib.org/Partitions_ZF.html

\import{set.tex}
\import{set/cons.tex}
\import{set/powerset.tex}

\section{Bipartitions}

\begin{abbreviation}\label{bipartitions}
    $C$ is partitioned by $A$ and $B$ iff $A,B\neq\emptyset$ and $A$ is disjoint from $B$ and $A\union B = C$.
\end{abbreviation}

\begin{definition}\label{bipartitions_of_a_set}
    $\bipartitions{X} = \{ p\in\pow{X}\times\pow{X}\mid \text{$X$ is partitioned by $\fst{p}$ and $\snd{p}$}\}$.
\end{definition}

\begin{abbreviation}\label{is_a_bipartition_of}
    $P$ is a bipartition of $X$ iff $P\in\bipartitions{X}$.
\end{abbreviation}

\begin{proposition}\label{bipartition_intro}
    Suppose $C$ is partitioned by $A$ and $B$.
    Then $(A, B)$ is a bipartition of $C$.
\end{proposition}
\begin{proof}
    $(A,B)\in\pow{C}\times\pow{C}$ by
        \auto{powerset_intro,times_tuple_intro,subseteq,union_iff}.
    $C$ is partitioned by $\fst{(A,B)}$ and $\snd{(A,B)}$
        by \auto{fst_eq,snd_eq}.
    Thus $(A, B)$ is a bipartition of $C$ by \obv{bipartitions_of_a_set}.
\end{proof}

\begin{proposition}\label{bipartition_elim}
    Suppose $(A, B)$ is a bipartition of $C$.
    Then $C$ is partitioned by $A$ and $B$.
\end{proposition}
\begin{proof}
    $\fst{(A, B)} = A$ by \obv{fst_eq}.
    $\snd{(A, B)} = B$ by \obv{snd_eq}.
    Follows by \obv{bipartitions_of_a_set}.
\end{proof}

\begin{proposition}\label{bipartitions_emptyset}
    $\bipartitions{\emptyset}$ is empty.
\end{proposition}
\begin{proof}
    We have for all $p\in\bipartitions{\emptyset}$
        we have $p\in\emptyset$
        by \auto{bipartitions_of_a_set,powerset_elim,subseteq_emptyset_iff,times_proj_elim}.
    Follows by \obv{emptyset}.
\end{proof}

\begin{proposition}\label{union_eq_cons_implies_union_remove_eq}
    Suppose $d\notin C$.
    Suppose $A\union B = \cons{d}{C}$.
    Suppose $A, B\neq \{d\}$.
    Then $\remove{d}{A}\union \remove{d}{B} = C$.
\end{proposition}
\begin{proof}
    We have for all $x\in \remove{d}{A}\union \remove{d}{B}$
        we have $x\in C$ by \auto{cons_iff,setminus,union_iff}.
    We have for all $x\in C$
        we have $x\in \remove{d}{A}\union \remove{d}{B}$
        by \auto{cons_iff,setminus,union_iff}.
    Follows by set extensionality.
\end{proof}

\begin{proposition}\label{bipartitions_cons}
    Suppose $d\notin C$.
    Suppose $\cons{d}{C}$ is partitioned by $A$ and $B$.
    Suppose $A, B\neq \{d\}$.
    Then $C$ is partitioned by $\remove{d}{A}$ and $\remove{d}{B}$.
\end{proposition}
\begin{proof}
    $\remove{d}{A},\remove{d}{B}\neq\emptyset$ by
        \auto{setminus,singleton_iff,setext}.
    $\remove{d}{A}$ is disjoint from $\remove{d}{B}$
        by \auto{disjoint,setminus}.
    $\remove{d}{A}\union \remove{d}{B} = C$ by \obv{union_eq_cons_implies_union_remove_eq}.
\end{proof}
